% =============================================================
%  Section 3 — Threat Model: Eight Attack Chains
% =============================================================

\section{Threat Model: Eight Attack Chains}\label{sec:threat}

The preceding section established that agent-to-agent (A2A) commerce
introduces transaction patterns for which no human behavioral prior
exists.  We now make this concrete: starting from the \emph{documented}
capabilities of the OpenClaw platform~\cite{openclaw2026}, we derive
eight attack chains that exploit the gap between what fraud-detection
systems assume and what autonomous agents can actually do.

Each chain is grounded in a specific API call or composition of calls
available to any OpenClaw session.  We emphasize that none of these
chains require zero-day exploits or undocumented behavior; they follow
directly from the platform's published interface.
\Cref{tab:attack_chains} summarizes all eight chains, their difficulty
rating, the invariant each violates, and the detection signal most
sensitive to the resulting anomaly.

% ---- Summary table ----
\begin{table}[t]
  \centering
  \caption{%
    Eight attack chains derived from OpenClaw platform capabilities.
    Difficulty reflects the sophistication required to execute the
    chain, not the difficulty of detection.
    \textsuperscript{*}Impossible under current human-assumption
    detection systems.
    \textsuperscript{\dag}Impossible under per-address scoring
    architectures.%
  }
  \label{tab:attack_chains}
  \small
  \begin{tabularx}{\textwidth}{@{}l l X l l@{}}
    \toprule
    \textbf{Chain} & \textbf{Attack} & \textbf{Invariant Violated}
      & \textbf{Difficulty} & \textbf{Primary Signal} \\
    \midrule
    CHAIN\_1 & Agent Enumeration
      & Account visibility     & Easy      & Network Topology \\
    CHAIN\_2 & History Extraction
      & Transaction privacy    & Medium    & NT, Value Flow \\
    CHAIN\_3 & Async Flooding
      & Velocity               & Medium    & Temporal Consistency \\
    CHAIN\_4 & Disposable Agent Army
      & Identity persistence   & Hard      & Network Topology \\
    CHAIN\_5 & Cross-Platform Identity
      & Channel isolation      & Impossible\textsuperscript{*}
      & Econ.\ Rationality, NT \\
    CHAIN\_6 & Behavioral Mimicry
      & Unique fingerprint     & Impossible\textsuperscript{*}
      & Temporal Consistency \\
    CHAIN\_7 & Swarm Intelligence
      & Coordination latency   & Impossible\textsuperscript{\dag}
      & Collective (gap) \\
    CHAIN\_8 & Market Manipulation
      & Reaction time          & Impossible\textsuperscript{*}
      & Value Flow, ER \\
    \bottomrule
  \end{tabularx}
\end{table}

The remainder of this section details the four chains that are most
instructive for detection design: async flooding (\cref{subsec:chain3}),
disposable agent armies (\cref{subsec:chain4}), the behavioral mimicry
paradox (\cref{subsec:chain6}), and swarm intelligence
(\cref{subsec:chain7}).  The other four chains are summarized in
\cref{subsec:chains_other}.


% ============================================================
\subsection{CHAIN\_3: Async Flooding}\label{subsec:chain3}

\paragraph{Mechanism.}
The OpenClaw \texttt{sessions\_send} API accepts a
\texttt{timeoutSeconds} parameter.  When set to zero, the call returns
immediately---it fires a transaction without waiting for confirmation.
An agent can therefore issue transactions as fast as the network will
accept them, bounded only by RPC round-trip latency (approximately
400\,ms on Base~L2).

\paragraph{Velocity arithmetic.}
The resulting throughput is straightforward to compute:
\begin{equation}\label{eq:async_velocity}
  v_{\text{async}} = \frac{1}{0.4\;\text{s}}
                    = 2.5\;\text{tx/s}
                    = 9{,}000\;\text{tx/h}.
\end{equation}
For comparison, a human operator interacting through a banking
interface sustains roughly 100 transactions per day on average,
or approximately $4.2$\;tx/h.  Even a power user executing
20~transactions per hour through a mobile app achieves at most
$v_{\text{human,peak}} \approx 20$\;tx/h.  The velocity
amplification factor comparing \emph{peak agent} to
\emph{peak human} throughput is:
\begin{equation}\label{eq:velocity_amp}
  \alpha_v = \frac{v_{\text{async}}}{v_{\text{human,peak}}}
           = \frac{9{,}000\;\text{tx/h}}{20\;\text{tx/h}}
           = 450.
\end{equation}
Against average human throughput the factor exceeds $10^{6}$;
against peak human throughput it remains ${\sim}450\times$---still
far beyond any velocity threshold in production rule engines.
Critically, this is not a theoretical upper bound---it is the
\emph{default} behavior of a single agent calling a documented
API with a single parameter change.

\paragraph{Why this matters for detection.}
No velocity cap in a traditional rule engine is calibrated for
throughput six orders of magnitude above human norms.
The Temporal Consistency signal (\cref{subsec:signals}) catches
async flooding through two sub-signals: burst detection, which
flags sustained inter-transaction intervals below 400\,ms, and
circadian violation, which flags activity outside any plausible
human waking window when sustained at this rate.  Importantly, an
attacker who throttles the flood to evade velocity rules
sacrifices the economic advantage of the attack, creating a
natural deterrent even without perfect detection.


% ============================================================
\subsection{CHAIN\_4: Disposable Agent Army}\label{subsec:chain4}

\paragraph{Mechanism.}
The \texttt{sessions\_spawn} API creates sub-agents with configurable
lifecycle parameters.  Two parameters are critical:
\texttt{cleanup: "delete"} removes all traces of the agent after the
session ends, and \texttt{pruneAfter: "1h"} sets an automatic
expiration.  Combined, these allow an orchestrator to create agents
that exist just long enough to execute a transaction and then vanish.

\paragraph{Scale.}
Each spawn call completes in under one second.  An orchestrator
can therefore instantiate 100 sub-agents in fewer than 100 seconds,
each of which receives a fresh address with a clean reputation
slate.  After one hour, every sub-agent is pruned, leaving no
persistent identity for reputation systems to score.

\paragraph{The identity-persistence assumption.}
Contemporary fraud-detection systems universally assume that an
identity, whether a bank account, a wallet address, or a user ID,
persists long enough to accumulate a behavioral history.  Reputation
scores, velocity counters, and behavioral models all depend on
this assumption.  The disposable agent army violates it
categorically: each agent's history is at most one hour long,
and typically consists of a single transaction.

\paragraph{Detection via Network Topology.}
Although individual disposable agents are difficult to fingerprint,
the \emph{orchestrator} that spawns them is not.  The Network
Topology signal (\cref{subsec:signals}) detects this pattern through
sender centrality: a single address that fans out to dozens of
previously unseen addresses within a narrow time window produces an
anomalous hub in the transaction graph.  The addresses themselves
may be ephemeral, but the star topology they form is a durable
structural invariant of the attack.


% ============================================================
\subsection{CHAIN\_6: Behavioral Mimicry Paradox}\label{subsec:chain6}

\paragraph{Mechanism.}
The \texttt{sessions\_history} API returns the complete behavioral
profile of any agent---its transaction timestamps, amounts,
counterparties, and message contents.  A malicious agent can query
this history, compute the target's statistical fingerprint (mean
inter-transaction interval, amount distribution, preferred
counterparties), and mimic it in its own transactions.

\paragraph{Mimicry quality.}
An agent programmed to replicate a target's timing distribution can
achieve a coefficient of variation (CV) in inter-transaction intervals
below 0.01---effectively perfect regularity.  This might seem like a
successful evasion: the mimic's transactions ``look like'' the
target's.

\paragraph{The paradox.}
In practice, this mimicry makes the agent \emph{more} detectable, not
less.  The reason is that real agents---even automated ones operating
legitimately---exhibit substantial timing variability due to network
latency jitter, queue contention, rate-limit backoffs, and varying
computational loads.  Empirically, we observe $\cvt = 1.87$ for
genuine agent traffic (see \cref{sec:evaluation}).  An entity
transacting with $\cvt < 0.05$ is performing an act that no real
system produces organically: it is \emph{too regular}.

\begin{equation}\label{eq:mimicry_paradox}
  \cvt_{\text{mimic}} < 0.05 \ll \cvt_{\text{real}} \approx 1.87
  \quad \Longrightarrow \quad
  \text{``too perfect'' regularity is an agent-invariant signal.}
\end{equation}

This is a fundamental result.  Traditional fraud detection looks for
anomalously \emph{irregular} behavior---spikes, outliers, deviations
from the norm.  Agent-invariant detection must also look in the
opposite direction: anomalous \emph{regularity} is equally
diagnostic.  The Temporal Consistency signal captures this through its
timing-regularity sub-signal, which fires on both tails of the CV
distribution (\cref{subsec:signals}).


% ============================================================
\subsection{CHAIN\_7: Swarm Intelligence}\label{subsec:chain7}

\paragraph{Mechanism.}
The \texttt{sessions\_send} API supports broadcast delivery: a single
message can instruct 100 or more agents simultaneously.  Each agent
in the swarm executes exactly one transaction---say, purchasing a
small amount of the same token---and then goes silent.

\paragraph{Individual invisibility.}
From the perspective of any per-address detection system, each swarm
member is unremarkable: a single transaction of modest value,
performed at a plausible time, from an address with no prior
suspicious history.  No velocity rule fires.  No behavioral model
has enough data to score.  The agent is individually invisible.

\paragraph{Collective visibility.}
Zoom out, however, and the swarm is unmistakable.  Fifty or more
addresses, previously dormant, execute transactions of similar
amounts to the same target within the same block or narrow time
window.  The coordination is visible in the population distribution,
even though it is invisible in any single address's history.

\paragraph{The architectural blind spot.}
This is not a detection failure per se---the signal is present in the
data.  It is an \emph{architectural} failure: current fraud-detection
systems score individual addresses, not populations.  There is no
module that asks ``did 50 unrelated addresses just do the same thing
at the same time?''  We identify this as a \textbf{collective blind
spot} in~\cref{tab:attack_chains}, and the absence of a validated
population-level signal is a known gap in our detection framework
(\cref{subsec:signals}).

\paragraph{Implications.}
Closing this gap requires a fundamentally different detection
primitive: one that operates on \emph{sets of addresses} rather than
individual addresses.  We discuss candidate approaches---including
graph-community detection over sliding windows and statistical tests
for coordinated activation---in \cref{sec:recommendations}, while
acknowledging that no validated implementation exists at the time of
writing.


% ============================================================
\subsection{Remaining Attack Chains}\label{subsec:chains_other}

\paragraph{CHAIN\_1: Agent Enumeration.}
The OpenClaw agent registry is publicly queryable.  An adversary can
enumerate all registered agents, their capabilities, and their
communication endpoints, building a complete map of the agent
ecosystem.  This violates the account-visibility invariant that
underpins privacy-by-obscurity assumptions in traditional finance.
The Network Topology signal detects enumeration through path-anomaly
scoring: systematic querying of many endpoints produces a distinctive
star-shaped access pattern.

\paragraph{CHAIN\_2: History Extraction.}
As noted in CHAIN\_6, the \texttt{sessions\_history} API exposes full
transaction histories.  Even without mimicry, this violates
transaction privacy: an adversary can reconstruct the complete
financial profile of any agent.  Detection relies on a combination of
Network Topology (unusual query patterns) and Value Flow (the
extracted data enabling subsequent layering attacks).

\paragraph{CHAIN\_5: Cross-Platform Identity.}
An agent operating on multiple blockchains simultaneously---Base,
Ethereum mainnet, Arbitrum---can exploit the channel-isolation
assumption: no single-chain fraud detector sees the full picture.
This chain is rated ``Impossible'' under current human-assumption
systems because humans rarely transact on multiple chains
simultaneously.  The Cross-Platform Correlation signal
(\cref{subsec:signals}) is designed for exactly this case, though it
remains inactive in our current single-chain implementation.

\paragraph{CHAIN\_8: Market Manipulation.}
Agents can react to on-chain events in milliseconds---faster than any
human trader and faster than most monitoring systems.  By front-running
price movements or exploiting stale oracle data, an agent can extract
value in ways that violate the reaction-time invariant of human-paced
markets.  The Value Flow and Economic Rationality signals detect the
resulting anomalies: transactions that are economically rational only
if the actor has sub-second information access.

\medskip
\noindent
Taken together, these eight chains demonstrate that the A2A commerce
surface is not merely larger than the human-commerce surface---it is
\emph{qualitatively different}.  The next section presents a detection
framework designed for this new surface.
